\documentclass[11pt,a4paper]{article}

\usepackage[utf8]{inputenc}
\usepackage{amsmath, amssymb, amsthm}
\usepackage{graphicx}
\usepackage[margin=1in]{geometry}
\usepackage[hidelinks]{hyperref}

% Numbered environments share the section counter, so Theorem 2.1 is
% the first theorem of Section 2 — easier to find while reading back.
\theoremstyle{plain}
\newtheorem{theorem}{Theorem}[section]
\newtheorem{lemma}[theorem]{Lemma}
\newtheorem{corollary}[theorem]{Corollary}

\theoremstyle{definition}
\newtheorem{definition}[theorem]{Definition}
\newtheorem{example}[theorem]{Example}

\theoremstyle{remark}
\newtheorem*{remark}{Remark}

\title{{name}}
\author{Your Name}
\date{\today}

\begin{document}

\maketitle

\tableofcontents

\section{Introduction}

What this set of notes covers, and what a reader is assumed to know
already.

\section{Definitions}

\begin{definition}[Continuity]
    \label{def:continuity}
    A function $f : \mathbb{R} \to \mathbb{R}$ is continuous at $a$ if
    for every $\varepsilon > 0$ there is a $\delta > 0$ such that
    $|x - a| < \delta$ implies $|f(x) - f(a)| < \varepsilon$.
\end{definition}

\begin{remark}
    Notes to yourself belong in a remark, where they are visually
    separate from the statements you will be examined on.
\end{remark}

\section{Results}

\begin{theorem}
    \label{thm:example}
    State the theorem precisely. Every symbol should already have been
    defined, as in Definition~\ref{def:continuity}.
\end{theorem}

\begin{proof}
    Give the argument. End it with the square, which \texttt{amsthm}
    places automatically. \qedhere
\end{proof}

\begin{corollary}
    A consequence of Theorem~\ref{thm:example} worth naming separately.
\end{corollary}

\begin{example}
    A concrete case. Notes without examples are hard to revise from.
\end{example}

\end{document}
